\documentclass[12pt]{article}

\usepackage{fontspec}
\newfontfamily\handwriting{a-casual-handwritten-pen}[Path=../assets/fonts/,Extension=.ttf,Scale=1.5]
\newfontfamily\comic{Comic Neue}

\usepackage[utf8]{inputenc}
\usepackage{amsmath}
\usepackage{amssymb}
\usepackage{enumitem}
\usepackage{xcolor}
\usepackage{soul}
\usepackage{tikz}
\usepackage[most]{tcolorbox}
\usepackage{titlesec}
\usepackage{multicol}
\usepackage{environ}

\usetikzlibrary{fit, positioning, arrows.meta, decorations.pathreplacing}

% --- Custom Colors ---
\definecolor{primaryblue}{RGB}{42, 111, 151}
\definecolor{exbg}{RGB}{255, 240, 245}
\definecolor{rulebg}{RGB}{232, 245, 233}
\definecolor{highlight}{RGB}{255, 243, 176}
\definecolor{pinkanno}{RGB}{255, 105, 180}
\definecolor{purpleanno}{RGB}{147, 112, 219}
\definecolor{solutionbg}{RGB}{245, 245, 255}

% --- Header Styling ---
\titleformat{\section}{\color{primaryblue}\normalfont\Large\bfseries}{\thesection}{1em}{}
\titleformat{\subsection}{\color{primaryblue}\normalfont\large\bfseries}{\thesubsection}{1em}{}

% --- Friendly Boxes ---
% --- Auto-numbered Example Box ---
\newcounter{examplenum}

\newtcolorbox{friendlyexbox}[1]{
    colback=exbg,
    colframe=pinkanno!50,
    arc=5pt,
    title=\textbf{#1},
    coltitle=black,
    attach title to upper,
    after title={:\quad}
}

\newenvironment{friendlyex}{%
    \refstepcounter{examplenum}%
    \begin{friendlyexbox}{Example \theexamplenum}%
}{%
    \end{friendlyexbox}%
}

\newtcolorbox{friendlyrule}[1]{
    colback=rulebg,
    colframe=primaryblue!30,
    arc=5pt,
    fonttitle=\bfseries,
    title={#1},
    coltitle=primaryblue
}

\newcommand{\funhl}[1]{\sethlcolor{highlight}\hl{#1}}

% --- Show/Hide Solutions ---
\newif\ifshowsolutions

% Uncomment the next line to show solutions.
\showsolutionstrue

\newcommand{\solutionblank}{\vspace{25ex}}

\NewEnviron{solutionbox}{
\ifshowsolutions
\begin{tcolorbox}[
    colback=solutionbg,
    colframe=purpleanno!45,
    arc=5pt,
    title={Solution},
    coltitle=primaryblue,
    fonttitle=\bfseries
]
\BODY
\end{tcolorbox}
\else
\solutionblank
\fi
}

\begin{document}
\handwriting

\begin{center}
    \Large\textbf{\color{primaryblue}Module 3 Lesson 2\\Graphs of Polynomial Functions} \\
    \vspace{0.2cm}
    \hrule
\end{center}

\section*{Objectives}

\begin{friendlyrule}{In this lesson, we will learn how to}
\begin{itemize}
    \item sketch the graph of a polynomial function using end behavior, intercepts, and multiplicity;
    \item read a polynomial's graph to recover its zeros, multiplicities, and a possible formula.
\end{itemize}
\end{friendlyrule}

\section*{1. Strategy for Sketching a Polynomial}

\begin{friendlyrule}{Steps to Sketch $f(x)$}
\begin{enumerate}
    \item Determine the \funhl{end behavior} (degree and sign of leading coefficient).
    \item Find the \funhl{$y$-intercept}: evaluate $f(0)$.
    \item Find the \funhl{zeros} (factor if necessary) and their multiplicities.
    \item Determine whether the graph \funhl{crosses} or \funhl{bounces} at each zero.
    \item Plot a few additional points if needed, then connect with a smooth, continuous curve consistent with the end behavior.
\end{enumerate}
\end{friendlyrule}

\begin{friendlyex}{}
Sketch the graph of $f(x)=-\dfrac{2}{3}(x+2)^2(x-3)$.
\end{friendlyex}

\begin{solutionbox}
\textbf{End behavior:} degree $3$ (odd), leading coefficient $-\frac23<0$: left up, right down.

\textbf{$y$-intercept:} $f(0)=-\frac23(2)^2(-3)=-\frac23(4)(-3)=8$, i.e.\ $(0,8)$.

\textbf{Zeros:} $x=-2$ (multiplicity $2$, bounce), $x=3$ (multiplicity $1$, cross).

\vspace{0.3cm}
\begin{center}
\begin{tikzpicture}
\begin{scope}[xscale=0.85,yscale=0.11]
\draw[-{Latex}] (-3.6,0) -- (4.2,0) node[right] {$x$};
\draw[-{Latex}] (0,-26) -- (0,15) node[above] {$y$};
\foreach \x in {-3,-2,-1,1,2,3,4} \draw (\x,-0.8) -- (\x,0.8) node[below=6pt] {\tiny $\x$};
\draw[teal, thick, domain=-3.35:3.75, samples=90, smooth] plot (\x, {-2/3*(\x+2)*(\x+2)*(\x-3)});
\filldraw[pinkanno] (-2,0) circle (2pt);
\filldraw[pinkanno] (3,0) circle (2pt);
\filldraw[pinkanno] (0,8) circle (2pt);
\node[anchor=south west] at (-3.5,10) {\small $(0,8)$};
\node[below] at (-2,-9) {\small $x=-2$};
\node[below] at (3,-9) {\small $x=3$};
\end{scope}
\end{tikzpicture}
\end{center}
\end{solutionbox}

\begin{friendlyex}{}
Sketch the graph of $f(x)=\dfrac{3}{8}(x+2)(x-1)^2(x-4)$.
\end{friendlyex}

\begin{solutionbox}
\textbf{End behavior:} degree $4$ (even), leading coefficient $\frac38>0$: both ends go up.

\textbf{$y$-intercept:} $f(0)=\frac38(2)(1)(-4)=\frac38(-8)=-3$, i.e.\ $(0,-3)$.

\textbf{Zeros:} $x=-2$ (multiplicity $1$, cross), $x=1$ (multiplicity $2$, bounce), $x=4$ (multiplicity $1$, cross).

\vspace{0.3cm}
\begin{center}
\begin{tikzpicture}
\begin{scope}[xscale=0.8,yscale=0.11]
\draw[-{Latex}] (-3.2,0) -- (5,0) node[right] {$x$};
\draw[-{Latex}] (0,-10) -- (0,42) node[above] {$y$};
\foreach \x in {-2,-1,1,2,3,4} \draw (\x,-1) -- (\x,1) node[below=6pt] {\tiny $\x$};
\draw[teal, thick, domain=-2.75:4.55, samples=90, smooth] plot (\x, {3/8*(\x+2)*(\x-1)*(\x-1)*(\x-4)});
\filldraw[pinkanno] (-2,0) circle (2pt);
\filldraw[pinkanno] (1,0) circle (2pt);
\filldraw[pinkanno] (4,0) circle (2pt);
\filldraw[pinkanno] (0,-3) circle (2pt);
\node[anchor=south east] at (0,4) {\small $(0,-3)$};
\node[below] at (-2,-4.5) {\small $x=-2$};
\node[above] at (1,4) {\small $x=1$};
\node[below] at (4,-4.5) {\small $x=4$};
\end{scope}
\end{tikzpicture}
\end{center}
\end{solutionbox}

\section*{2. Reading a Graph}

\begin{friendlyrule}{Reading Zeros and Multiplicity from a Graph}
\begin{itemize}
    \item Where the graph \funhl{crosses} the $x$-axis, the corresponding zero has \funhl{odd} multiplicity (often $1$).
    \item Where the graph \funhl{bounces} (touches and turns), the corresponding zero has \funhl{even} multiplicity (often $2$).
    \item Multiply the leading coefficient's sign and the end behavior together with the zeros to build a possible formula $f(x)=a(x-r_1)^{m_1}(x-r_2)^{m_2}\cdots$.
\end{itemize}
\end{friendlyrule}

\begin{friendlyex}{}
The graph of a degree-$4$ polynomial is shown below. Determine its zeros, their multiplicities, and a possible formula for $f(x)$.

\begin{center}
\begin{tikzpicture}
\begin{scope}[xscale=0.95,yscale=0.5]
\draw[-{Latex}] (-4,0) -- (3.5,0) node[right] {$x$};
\draw[-{Latex}] (0,-5) -- (0,10) node[above] {$y$};
\foreach \x in {-3,-2,-1,1,2,3} \draw (\x,-0.2) -- (\x,0.2) node[below=4pt] {\tiny $\x$};
\draw[thick, domain=-3.55:3.15, samples=90, smooth] plot (\x, {1/4*(\x+3)*(\x+1)*(\x-2)*(\x-2)});
\filldraw[pinkanno] (-3,0) circle (2pt);
\filldraw[pinkanno] (-1,0) circle (2pt);
\filldraw[pinkanno] (2,0) circle (2pt);
\filldraw[pinkanno] (0,3) circle (2pt);
\node[above] at (0.6,3.3) {\small $(0,3)$};
\end{scope}
\end{tikzpicture}
\end{center}
\end{friendlyex}

\begin{solutionbox}
The graph \textbf{crosses} at $x=-3$ and $x=-1$ (odd multiplicity, take $1$), and \textbf{bounces} at $x=2$ (even multiplicity, take $2$). This already accounts for degree $1+1+2=4$, matching the given degree.

A possible formula:
\[
f(x)=a(x+3)(x+1)(x-2)^2.
\]

Use the $y$-intercept $(0,3)$ to find $a$:
\[
3=a(3)(1)(4)=12a \implies a=\frac14.
\]

\[
\boxed{f(x)=\frac14(x+3)(x+1)(x-2)^2}
\]
\end{solutionbox}

\begin{friendlyex}{}
The graph of a degree-$5$ polynomial is shown below. Determine its zeros, their multiplicities, and a possible formula for $f(x)$.

\begin{center}
\begin{tikzpicture}
\begin{scope}[xscale=1.1,yscale=0.4]
\draw[-{Latex}] (-3.2,0) -- (2.6,0) node[right] {$x$};
\draw[-{Latex}] (0,-5.5) -- (0,7) node[above] {$y$};
\foreach \x in {-2,-1,1,2} \draw (\x,-0.25) -- (\x,0.25) node[below=4pt] {\tiny $\x$};
\draw[thick, domain=-2.4:1.9, samples=90, smooth] plot (\x, {1/2*(\x+2)*(\x+2)*(\x-1)*(\x-1)*(\x-1)});
\filldraw[pinkanno] (-2,0) circle (2pt);
\filldraw[pinkanno] (1,0) circle (2pt);
\filldraw[pinkanno] (0,-2) circle (2pt);
\node[below] at (0.55,-2.4) {\small $(0,-2)$};
\end{scope}
\end{tikzpicture}
\end{center}
\end{friendlyex}

\begin{solutionbox}
The graph \textbf{bounces} at $x=-2$ (even multiplicity, take $2$) and \textbf{crosses} with a flattened S-shape at $x=1$ (odd multiplicity greater than $1$, take $3$). This accounts for degree $2+3=5$, matching the given degree.

A possible formula:
\[
f(x)=a(x+2)^2(x-1)^3.
\]

Use the $y$-intercept $(0,-2)$ to find $a$:
\[
-2=a(4)(-1)=-4a \implies a=\frac12.
\]

\[
\boxed{f(x)=\frac12(x+2)^2(x-1)^3}
\]
\end{solutionbox}

\end{document}
